% Options for packages loaded elsewhere
\PassOptionsToPackage{unicode}{hyperref}
\PassOptionsToPackage{hyphens}{url}
%
\documentclass[
  brazilian,
]{article}
\usepackage{lmodern}
\usepackage{amssymb,amsmath}
\usepackage{ifxetex,ifluatex}
\ifnum 0\ifxetex 1\fi\ifluatex 1\fi=0 % if pdftex
  \usepackage[T1]{fontenc}
  \usepackage[utf8]{inputenc}
  \usepackage{textcomp} % provide euro and other symbols
\else % if luatex or xetex
  \usepackage{unicode-math}
  \defaultfontfeatures{Scale=MatchLowercase}
  \defaultfontfeatures[\rmfamily]{Ligatures=TeX,Scale=1}
\fi
% Use upquote if available, for straight quotes in verbatim environments
\IfFileExists{upquote.sty}{\usepackage{upquote}}{}
\IfFileExists{microtype.sty}{% use microtype if available
  \usepackage[]{microtype}
  \UseMicrotypeSet[protrusion]{basicmath} % disable protrusion for tt fonts
}{}
\makeatletter
\@ifundefined{KOMAClassName}{% if non-KOMA class
  \IfFileExists{parskip.sty}{%
    \usepackage{parskip}
  }{% else
    \setlength{\parindent}{0pt}
    \setlength{\parskip}{6pt plus 2pt minus 1pt}}
}{% if KOMA class
  \KOMAoptions{parskip=half}}
\makeatother
\usepackage{xcolor}
\IfFileExists{xurl.sty}{\usepackage{xurl}}{} % add URL line breaks if available
\IfFileExists{bookmark.sty}{\usepackage{bookmark}}{\usepackage{hyperref}}
\hypersetup{
  pdftitle={Parte III: Processamento de Linguagem Natural para Linguistas},
  pdfauthor={CiberExt 26-29 · FEELT38103 · Universidade Federal de Uberlândia},
  pdflang={pt-BR},
  hidelinks,
  pdfcreator={LaTeX via pandoc}}
\urlstyle{same} % disable monospaced font for URLs
\usepackage{color}
\usepackage{fancyvrb}
\newcommand{\VerbBar}{|}
\newcommand{\VERB}{\Verb[commandchars=\\\{\}]}
\DefineVerbatimEnvironment{Highlighting}{Verbatim}{commandchars=\\\{\}}
% Add ',fontsize=\small' for more characters per line
\newenvironment{Shaded}{}{}
\newcommand{\AlertTok}[1]{\textcolor[rgb]{1.00,0.00,0.00}{\textbf{#1}}}
\newcommand{\AnnotationTok}[1]{\textcolor[rgb]{0.38,0.63,0.69}{\textbf{\textit{#1}}}}
\newcommand{\AttributeTok}[1]{\textcolor[rgb]{0.49,0.56,0.16}{#1}}
\newcommand{\BaseNTok}[1]{\textcolor[rgb]{0.25,0.63,0.44}{#1}}
\newcommand{\BuiltInTok}[1]{#1}
\newcommand{\CharTok}[1]{\textcolor[rgb]{0.25,0.44,0.63}{#1}}
\newcommand{\CommentTok}[1]{\textcolor[rgb]{0.38,0.63,0.69}{\textit{#1}}}
\newcommand{\CommentVarTok}[1]{\textcolor[rgb]{0.38,0.63,0.69}{\textbf{\textit{#1}}}}
\newcommand{\ConstantTok}[1]{\textcolor[rgb]{0.53,0.00,0.00}{#1}}
\newcommand{\ControlFlowTok}[1]{\textcolor[rgb]{0.00,0.44,0.13}{\textbf{#1}}}
\newcommand{\DataTypeTok}[1]{\textcolor[rgb]{0.56,0.13,0.00}{#1}}
\newcommand{\DecValTok}[1]{\textcolor[rgb]{0.25,0.63,0.44}{#1}}
\newcommand{\DocumentationTok}[1]{\textcolor[rgb]{0.73,0.13,0.13}{\textit{#1}}}
\newcommand{\ErrorTok}[1]{\textcolor[rgb]{1.00,0.00,0.00}{\textbf{#1}}}
\newcommand{\ExtensionTok}[1]{#1}
\newcommand{\FloatTok}[1]{\textcolor[rgb]{0.25,0.63,0.44}{#1}}
\newcommand{\FunctionTok}[1]{\textcolor[rgb]{0.02,0.16,0.49}{#1}}
\newcommand{\ImportTok}[1]{#1}
\newcommand{\InformationTok}[1]{\textcolor[rgb]{0.38,0.63,0.69}{\textbf{\textit{#1}}}}
\newcommand{\KeywordTok}[1]{\textcolor[rgb]{0.00,0.44,0.13}{\textbf{#1}}}
\newcommand{\NormalTok}[1]{#1}
\newcommand{\OperatorTok}[1]{\textcolor[rgb]{0.40,0.40,0.40}{#1}}
\newcommand{\OtherTok}[1]{\textcolor[rgb]{0.00,0.44,0.13}{#1}}
\newcommand{\PreprocessorTok}[1]{\textcolor[rgb]{0.74,0.48,0.00}{#1}}
\newcommand{\RegionMarkerTok}[1]{#1}
\newcommand{\SpecialCharTok}[1]{\textcolor[rgb]{0.25,0.44,0.63}{#1}}
\newcommand{\SpecialStringTok}[1]{\textcolor[rgb]{0.73,0.40,0.53}{#1}}
\newcommand{\StringTok}[1]{\textcolor[rgb]{0.25,0.44,0.63}{#1}}
\newcommand{\VariableTok}[1]{\textcolor[rgb]{0.10,0.09,0.49}{#1}}
\newcommand{\VerbatimStringTok}[1]{\textcolor[rgb]{0.25,0.44,0.63}{#1}}
\newcommand{\WarningTok}[1]{\textcolor[rgb]{0.38,0.63,0.69}{\textbf{\textit{#1}}}}
\setlength{\emergencystretch}{3em} % prevent overfull lines
\providecommand{\tightlist}{%
  \setlength{\itemsep}{0pt}\setlength{\parskip}{0pt}}
\setcounter{secnumdepth}{-\maxdimen} % remove section numbering
\ifxetex
  % Load polyglossia as late as possible: uses bidi with RTL langages (e.g. Hebrew, Arabic)
  \usepackage{polyglossia}
  \setmainlanguage[]{brazilian}
\else
  \usepackage[shorthands=off,main=brazilian]{babel}
\fi

\title{Parte III: Processamento de Linguagem Natural para Linguistas}
\usepackage{etoolbox}
\makeatletter
\providecommand{\subtitle}[1]{% add subtitle to \maketitle
  \apptocmd{\@title}{\par {\large #1 \par}}{}{}
}
\makeatother
\subtitle{Unidade 6 - Anotações de nível de discurso}
\author{CiberExt 26-29 · FEELT38103 · Universidade Federal de
Uberlândia}
\date{Agosto de 2026}

\begin{document}
\maketitle

\hypertarget{trabalhando-com-anotauxe7uxf5es-de-nuxedvel-de-discurso}{%
\section{Trabalhando com anotações de nível de
discurso}\label{trabalhando-com-anotauxe7uxf5es-de-nuxedvel-de-discurso}}

Nas seções anteriores, exploramos diversas técnicas de Processamento de
Linguagem Natural (PLN) para gerar anotações linguísticas a partir de
textos puros.

Agora, o nosso foco será aprender a aproveitar anotações linguísticas
que já estão prontas, dando atenção especial àquelas que analisam
estruturas maiores do que uma simples oração (cláusula) --- os chamados
fenômenos de nível de discurso.

Ao final desta seção, você deverá ser capaz de:

\begin{itemize}
\tightlist
\item
  Compreender o básico do formato de anotação CoNLL-U.
\item
  Criar um objeto \emph{Doc} do spaCy manualmente, a partir do zero.
\item
  Anotar trechos de texto (\emph{Spans}) dentro de um objeto \emph{Doc}
  utilizando \emph{SpanGroups}.
\item
  Importar corpora com anotações no formato CoNLL-U para o ecossistema
  do spaCy.
\end{itemize}

\hypertarget{o-formato-de-anotauxe7uxe3o-conll-u}{%
\subsection{O formato de anotação
CoNLL-U}\label{o-formato-de-anotauxe7uxe3o-conll-u}}

O CoNLL-X é um formato padrão para descrever características
linguísticas em uma grande variedade de idiomas (Buchholz e Marsi
\href{https://www.aclweb.org/anthology/W06-2920}{2006}). Ele foi criado,
a princípio, para facilitar o compartilhamento de dados nas chamadas
\emph{shared tasks} (desafios computacionais abertos a pesquisadores) no
campo do PLN (veja Nissim et
al.~\href{https://doi.org/10.1162/COLI_a_00304}{2017}).

O \href{https://universaldependencies.org/format.html}{CoNLL-U} é a
evolução desse formato, adaptado para suportar o modelo de Dependências
Universais (\emph{Universal Dependencies}), do qual falamos na Parte
III. Hoje em dia, o CoNLL-U é amplamente usado para distribuir corpora
textuais em projetos baseados no \emph{Universal Dependencies}, mas
também é o queridinho de vários outros projetos independentes de
linguística computacional.

Para se ter uma ideia da sua abrangência, existem arquivos CoNLL-U não
só para \href{https://universaldependencies.org/}{inúmeras línguas
modernas}, mas até mesmo para idiomas da antiguidade, como o Acádio
(Luukko et
al.~\href{https://www.aclweb.org/anthology/2020.tlt-1.11}{2020}) e o
Copta (Zeldes e Abrams
\href{https://www.aclweb.org/anthology/W18-6022}{2018}).

\hypertarget{como-o-conll-u-funciona-por-dentro}{%
\subsubsection{Como o CoNLL-U funciona por
dentro}\label{como-o-conll-u-funciona-por-dentro}}

Os arquivos com as anotações CoNLL-U são sempre distribuídos como texto
simples (texto puro, como vimos na Parte II).

Ao abrir um desses arquivos, você notará três tipos de linhas:
\textbf{linhas de comentário} (comment lines), \textbf{linhas de
palavra} (word lines) e \textbf{linhas em branco} (blank lines).

\begin{itemize}
\tightlist
\item
  As \textbf{linhas de comentário} ficam antes do texto principal e
  começam com o símbolo da cerquilha (\texttt{\#}). Elas guardam
  informações extras (metadados) sobre o trecho de texto que vem a
  seguir.
\item
  As \textbf{linhas de palavra} carregam todas as anotações referentes a
  uma única palavra ou pontuação (token). As frases são formadas por uma
  sequência contínua dessas linhas de palavra.
\end{itemize}

A estrutura dessas anotações obedece a uma organização rigorosa de
colunas separadas pela tecla \emph{Tab}, contendo 10 campos essenciais:

\begin{verbatim}
ID  FORM    LEMMA   UPOS    XPOS    FEATS   HEAD    DEPREL  DEPS    MISC
\end{verbatim}

\begin{enumerate}
\def\labelenumi{\arabic{enumi}.}
\tightlist
\item
  \texttt{ID}: A posição (índice) da palavra na frase.
\item
  \texttt{FORM}: Como a palavra (ou pontuação) aparece no texto
  original.
\item
  \texttt{LEMMA}: A forma base ou de dicionário da palavra (lema).
\item
  \texttt{UPOS}: A
  \href{https://universaldependencies.org/u/pos/}{categoria gramatical
  universal} (substantivo, verbo, etc.).
\item
  \texttt{XPOS}: Uma categoria gramatical mais detalhada, específica do
  idioma analisado.
\item
  \texttt{FEATS}:
  \href{https://universaldependencies.org/u/feat/index.html}{Atributos
  morfológicos} (gênero, número, tempo verbal, etc.).
\item
  \texttt{HEAD}: O ID da palavra que atua como o ``chefe'' sintático
  (núcleo) da palavra atual.
\item
  \texttt{DEPREL}: O tipo de relação de dependência sintática que a
  palavra tem com seu \texttt{HEAD}.
\item
  \texttt{DEPS}:
  \href{https://universaldependencies.org/u/overview/enhanced-syntax.html}{Relações
  de dependência aprimoradas}.
\item
  \texttt{MISC}: Informações adicionais de qualquer natureza (um campo
  coringa).
\end{enumerate}

Por fim, uma \textbf{linha em branco} aparece no final das linhas de
palavras para sinalizar ao sistema que a frase terminou e outra vai
começar.

\hypertarget{trabalhando-com-conll-u-em-python}{%
\subsubsection{Trabalhando com CoNLL-U em
Python}\label{trabalhando-com-conll-u-em-python}}

Para manipular esses arquivos de forma inteligente usando o Python, a
nossa melhor amiga será a biblioteca
\href{https://github.com/EmilStenstrom/conllu/}{conllu}. Ela é leve e
muito competente para ler textos no formato CoNLL-U e transformá-los em
estruturas nativas do Python, fáceis de processar.

\begin{Shaded}
\begin{Highlighting}[]
\CommentTok{\# Importa a biblioteca conllu}
\ImportTok{import}\NormalTok{ conllu}
\end{Highlighting}
\end{Shaded}

Para o nosso estudo, vamos usar um arquivo CoNLL-U extraído do
\href{https://corpling.uis.georgetown.edu/gum/}{Georgetown University
Multilayer Corpus} (GUM; veja Zeldes
\href{http://dx.doi.org/10.1007/s10579-016-9343-x}{2017}). Vamos abrir o
arquivo como texto puro, lê-lo com o método \texttt{read()} e salvar o
texto na variável \texttt{annotations}.

\begin{Shaded}
\begin{Highlighting}[]
\CommentTok{\# Abre o arquivo de texto puro para leitura; atribui a \textquotesingle{}data\textquotesingle{}}
\ControlFlowTok{with} \BuiltInTok{open}\NormalTok{(}\StringTok{\textquotesingle{}data/GUM\_whow\_parachute.conllu\textquotesingle{}}\NormalTok{, mode}\OperatorTok{=}\StringTok{"r"}\NormalTok{, encoding}\OperatorTok{=}\StringTok{"utf{-}8"}\NormalTok{) }\ImportTok{as}\NormalTok{ data:}
    
    \CommentTok{\# Lê o conteúdo do arquivo e atribui a \textquotesingle{}annotations\textquotesingle{}}
\NormalTok{    annotations }\OperatorTok{=}\NormalTok{ data.read()}

\CommentTok{\# Verifica o tipo do objeto resultante}
\BuiltInTok{type}\NormalTok{(annotations)}
\end{Highlighting}
\end{Shaded}

\begin{verbatim}
str
\end{verbatim}

O arquivo foi carregado na memória do computador como uma gigantesca
string (texto contínuo). Podemos espiar os primeiros 1000 caracteres
dela para ter uma noção:

\begin{Shaded}
\begin{Highlighting}[]
\CommentTok{\# Imprime os 1000 primeiros caracteres da string em \textquotesingle{}annotations\textquotesingle{}}
\BuiltInTok{print}\NormalTok{(annotations[:}\DecValTok{1000}\NormalTok{])}
\end{Highlighting}
\end{Shaded}

\begin{verbatim}
# newdoc id = GUM_whow_parachute
# sent_id = GUM_whow_parachute-1
# text = How to Cope With a Double Parachute Failure
# s_type = inf
1   How how SCONJ   WRB PronType=Int    3   mark    _   Discourse=preparation:1->11
2   to  to  PART    TO  _   3   mark    _   _
3   Cope    Cope    VERB    VB  VerbForm=Inf    0   root    _   _
4   With    with    ADP IN  _   8   case    _   _
5   a   a   DET DT  Definite=Ind|PronType=Art   8   det _   Entity=(event-1
6   Double  double  ADJ JJ  Degree=Pos  8   amod    _   _
7   Parachute   parachute   NOUN    NN  Number=Sing 8   compound    _   Entity=(object-2)
8   Failure failure NOUN    NN  Number=Sing 3   obl _   Entity=event-1)

# sent_id = GUM_whow_parachute-2
# text = While skydiving, it is possible (yet extremely unlikely) that both your primary and reserve parachutes will malfunction, leaving you with no method of reducing your velocity.
# s_type = decl
1   While   while   SCONJ   IN  _   2   mark    _   Discourse=circumstance:2->3
2   skydiving   skydiving   NOUN    NN  Number=Sing 6   advcl   _   Entity=(event-3)|SpaceAfter=No
3   ,   ,   PUNCT   ,   _   2   punct   _   _
4   it  it  PRON    PRP Case=Nom|Gender=Neut|Num
\end{verbatim}

Fica bem nítida a estrutura que descrevemos acima: as linhas comentadas
com a cerquilha (\texttt{\#}), logo seguidas por colunas super
organizadas com os dados da palavra e suas análises linguísticas.

O símbolo sublinhado (\texttt{\_}) é usado quando uma das 10 colunas não
tem uma anotação definida (valor ausente).

Observe a décima coluna, a \texttt{MISC}. O corpus GUM costuma injetar
campos como \texttt{Discourse} e \texttt{Entity} para marcar informações
de nível superior sobre a estrutura do discurso e das entidades no texto
(quem é quem na história, lugares, objetos, etc.).

A grande questão agora é: como automatizar a leitura disso no Python sem
nos perdermos nessa string gigante?

Aqui entra a função \texttt{parse()} do módulo \texttt{conllu}. Ela faz
exatamente o trabalho duro de ``fatiar e traduzir'' a string bruta em um
formato legível.

\begin{Shaded}
\begin{Highlighting}[]
\CommentTok{\# Usa a função parse() para analisar as anotações; armazena em \textquotesingle{}sentences\textquotesingle{}}
\NormalTok{sentences }\OperatorTok{=}\NormalTok{ conllu.parse(annotations)}
\end{Highlighting}
\end{Shaded}

O que a função \texttt{parse()} nos devolve é uma lista do Python
recheada de objetos \emph{TokenList}. O \emph{TokenList} é a estrutura
de dados chave criada pela própria biblioteca \texttt{conllu} para
abrigar orações.

Vamos verificar o que se esconde no primeiro item
(\texttt{sentences{[}0{]}}).

\begin{Shaded}
\begin{Highlighting}[]
\NormalTok{sentences[}\DecValTok{0}\NormalTok{]}
\end{Highlighting}
\end{Shaded}

\begin{verbatim}
TokenList<How, to, Cope, With, a, Double, Parachute, Failure, metadata={newdoc id: "GUM_whow_parachute", sent_id: "GUM_whow_parachute-1", text: "How to Cope With a Double Parachute Failure", s_type: "inf"}>
\end{verbatim}

Maravilha, a biblioteca capturou os metadados (que estavam com
\texttt{\#}) e os disponibilizou sob a propriedade \texttt{metadata}.
Vamos acessar esse dicionário de metadados diretamente:

\begin{Shaded}
\begin{Highlighting}[]
\CommentTok{\# Obtém os metadados do primeiro item da lista}
\NormalTok{sentences[}\DecValTok{0}\NormalTok{].metadata}
\end{Highlighting}
\end{Shaded}

\begin{verbatim}
{'newdoc id': 'GUM_whow_parachute',
 'sent_id': 'GUM_whow_parachute-1',
 'text': 'How to Cope With a Double Parachute Failure',
 's_type': 'inf'}
\end{verbatim}

Isso nos revela como o projeto GUM costuma estruturar a papelada: ele
usa \texttt{newdoc\_id} para rotular o documento inteiro,
\texttt{sent\_id} para rotular cada sentença isolada, \texttt{text} com
a frase original limpa, e \texttt{s\_type} que categoriza a sentença sob
um viés gramatical (afirmação, pergunta, imperativo etc). (Zeldes \&
Simonson \href{https://www.aclweb.org/anthology/W16-1709}{2017}: 69).

Embora o terminal retorne um formato que parece um simples dicionário, a
estrutura por trás do \texttt{metadata} é um objeto especializado
(chamado \emph{Metadata}). No entanto, sua operação diária é idêntica a
um dicionário nativo do Python: bastam chaves e valores.

Portanto, se precisarmos apenas do modo verbal da oração
(\texttt{s\_type}), usamos a sintaxe clássica de dicionário:

\begin{Shaded}
\begin{Highlighting}[]
\CommentTok{\# Obtém o tipo de sentença sob a chave \textquotesingle{}s\_type\textquotesingle{}}
\NormalTok{sentences[}\DecValTok{0}\NormalTok{].metadata[}\StringTok{\textquotesingle{}s\_type\textquotesingle{}}\NormalTok{]}
\end{Highlighting}
\end{Shaded}

\begin{verbatim}
'inf'
\end{verbatim}

Isso nos retornará a string \texttt{inf}, indicando que a frase está no
modo infinitivo.

E sobre as palavras individuais? Como o próprio nome aponta, o objeto
\emph{TokenList} é construído agrupando múltiplos \emph{Tokens}.

Acessar o primeiro \emph{Token} do nosso primeiro \emph{TokenList} é
simples:

\begin{Shaded}
\begin{Highlighting}[]
\CommentTok{\# Obtém o primeiro token da primeira sentença}
\NormalTok{sentences[}\DecValTok{0}\NormalTok{][}\DecValTok{0}\NormalTok{]}
\end{Highlighting}
\end{Shaded}

\begin{verbatim}
{'id': 1,
 'form': 'How',
 'lemma': 'how',
 'upos': 'SCONJ',
 'xpos': 'WRB',
 'feats': {'PronType': 'Int'},
 'head': 3,
 'deprel': 'mark',
 'deps': None,
 'misc': {'Discourse': 'preparation:1->11'}}
\end{verbatim}

Mais uma vez, recebemos as informações num arranjo que emula a estrutura
de um dicionário.

Note bem a chave \texttt{misc}. O dicionário embutido nela carrega
anotações cruciais sobre as ``relações de discurso''. Essas relações
servem para amarrar ideias através de uma teoria gramatical profunda,
chamada \href{https://www.sfu.ca/rst}{Teoria da Estrutura Retórica (RST
- Rhetorical Structure Theory)} (Mann \& Thompson
\href{https://doi.org/10.1515/text.1.1988.8.3.243}{1988}).

De forma resumida, essa anotação nos diz que existe uma relação
discursiva (do tipo \textbf{preparation}) ligando as partes ``1'' e
``11'' do texto.

Porém, atenção: esses números não apontam para índices de palavras ou de
sentenças soltas. Eles apontam para ``Unidades Elementares do Discurso''
(do inglês \emph{Elementary Discourse Units}, ou EDU).

A segmentação em EDUs cria uma lente de aumento na qual os analistas
subdividem a narrativa em fragmentos básicos e avaliam de que forma eles
formam o esqueleto da história.

No molde clássico da Teoria da Estrutura Retórica (RST), as EDUs
coincidem majoritariamente com cláusulas ou orações, mas não se trata de
uma regra universal matemática.

\hypertarget{levando-a-anuxe1lise-de-discurso-para-objetos-doc-do-spacy}{%
\subsection{\texorpdfstring{Levando a análise de discurso para objetos
\emph{Doc} do
spaCy}{Levando a análise de discurso para objetos Doc do spaCy}}\label{levando-a-anuxe1lise-de-discurso-para-objetos-doc-do-spacy}}

Como verificamos anteriormente, os metadados embutidos no campo
\texttt{misc} (como a anotação \texttt{Discourse}) sinalizam o exato
momento onde uma nova EDU \textbf{inicia}.

Se percorrermos a nossa sequência de \emph{Tokens} rastreando o
surgimento dessas anotações, seremos capazes de desenhar precisamente as
fronteiras entre cada Unidade Elementar do Discurso --- que muitas vezes
cruzam o limite fixo das orações normais da nossa \emph{TokenList}.

A estratégia prática é varrer todos os \emph{Tokens} criando uma
contagem progressiva e anotar qual índice abriga o valioso carimbo
\texttt{Discourse}.

Vamos instanciar a variável de contagem (\texttt{counter}) e preparar o
``terreno'' (várias listas) para guardar com segurança as descobertas:

\begin{Shaded}
\begin{Highlighting}[]
\CommentTok{\# Cria uma variável com valor 0, que usaremos para contar}
\CommentTok{\# os Tokens processados}
\NormalTok{counter }\OperatorTok{=} \DecValTok{0}

\CommentTok{\# Usamos estas listas para acompanhar as sentenças, as unidades de discurso}
\CommentTok{\# e as relações que se mantêm entre elas.}
\NormalTok{discourse\_units }\OperatorTok{=}\NormalTok{ []}
\NormalTok{sent\_types }\OperatorTok{=}\NormalTok{ []}
\NormalTok{relations }\OperatorTok{=}\NormalTok{ []}

\CommentTok{\# Cria listas vazias para a informação que extrairemos das anotações}
\CommentTok{\# CoNLL{-}U. Essas listas serão usadas para criar um objeto Doc do}
\CommentTok{\# spaCy mais adiante.}
\NormalTok{words }\OperatorTok{=}\NormalTok{ []}
\NormalTok{spaces }\OperatorTok{=}\NormalTok{ []}
\NormalTok{sent\_starts }\OperatorTok{=}\NormalTok{ []}
\end{Highlighting}
\end{Shaded}

O bloco a seguir irá analisar, palavra por palavra de cada
\texttt{sentence}, preenchendo as listas acima com todos os dados
relevantes (marcando também onde iniciam os EDUs):

\begin{Shaded}
\begin{Highlighting}[]
\CommentTok{\# Percorre cada objeto TokenList}
\ControlFlowTok{for}\NormalTok{ sentence }\KeywordTok{in}\NormalTok{ sentences:}
    
    \CommentTok{\# Ao começar a percorrer uma nova sentença, define o valor da}
    \CommentTok{\# variável \textquotesingle{}is\_start\textquotesingle{} como True. Isso marca o início de uma}
    \CommentTok{\# sentença.}
\NormalTok{    is\_start }\OperatorTok{=} \VariableTok{True}
    
    \CommentTok{\# Acrescenta o tipo da sentença atual à lista \textquotesingle{}sent\_types\textquotesingle{}}
\NormalTok{    sent\_types.append(sentence.metadata[}\StringTok{\textquotesingle{}s\_type\textquotesingle{}}\NormalTok{])}
        
    \CommentTok{\# Prossegue percorrendo os Tokens do TokenList atual}
    \ControlFlowTok{for}\NormalTok{ token }\KeywordTok{in}\NormalTok{ sentence:}
        
        \CommentTok{\# Usa a chave \textquotesingle{}form\textquotesingle{} para recuperar a forma da palavra do Token}
        \CommentTok{\# e a acrescenta à lista \textquotesingle{}words\textquotesingle{}.}
\NormalTok{        words.append(token[}\StringTok{\textquotesingle{}form\textquotesingle{}}\NormalTok{])}
        
        \CommentTok{\# Verifica se este Token inicia uma sentença, avaliando se a}
        \CommentTok{\# variável \textquotesingle{}is\_start\textquotesingle{} é True.}
        \ControlFlowTok{if}\NormalTok{ is\_start:}
            
            \CommentTok{\# Se o Token inicia uma sentença, acrescenta o valor True à lista}
            \CommentTok{\# \textquotesingle{}sent\_starts\textquotesingle{}. Note a ausência de aspas: trata{-}se de um valor}
            \CommentTok{\# booleano (True / False).}
\NormalTok{            sent\_starts.append(}\VariableTok{True}\NormalTok{)}
            
            \CommentTok{\# Define \textquotesingle{}is\_start\textquotesingle{} como False até que a próxima sentença comece}
            \CommentTok{\# e a variável seja novamente definida como True.}
\NormalTok{            is\_start }\OperatorTok{=} \VariableTok{False}
        
        \CommentTok{\# Se a variável \textquotesingle{}is\_start\textquotesingle{} for False, executa o bloco abaixo}
        \ControlFlowTok{else}\NormalTok{:}
            
            \CommentTok{\# Acrescenta o valor False à lista \textquotesingle{}sent\_starts\textquotesingle{}}
\NormalTok{            sent\_starts.append(}\VariableTok{False}\NormalTok{)}
        
        \CommentTok{\# Verifica se a chave \textquotesingle{}misc\textquotesingle{} contém algo e se ela guarda o valor}
        \CommentTok{\# \textquotesingle{}Discourse\textquotesingle{}; em caso afirmativo, executa o bloco abaixo.}
        \ControlFlowTok{if}\NormalTok{ token[}\StringTok{\textquotesingle{}misc\textquotesingle{}}\NormalTok{] }\KeywordTok{is} \KeywordTok{not} \VariableTok{None} \KeywordTok{and} \StringTok{\textquotesingle{}Discourse\textquotesingle{}} \KeywordTok{in}\NormalTok{ token[}\StringTok{\textquotesingle{}misc\textquotesingle{}}\NormalTok{]:}
            
            \CommentTok{\# A presença da chave \textquotesingle{}Discourse\textquotesingle{} indica o início de uma nova}
            \CommentTok{\# unidade elementar de discurso; acrescenta seu índice à lista}
            \CommentTok{\# \textquotesingle{}discourse\_units\textquotesingle{}.}
\NormalTok{            discourse\_units.append(counter)}
            
            \CommentTok{\# Desempacota a definição da relação; começa separando o nome da}
            \CommentTok{\# relação das unidades elementares de discurso. Atribui os objetos}
            \CommentTok{\# resultantes a \textquotesingle{}relation\textquotesingle{} e \textquotesingle{}edus\textquotesingle{}.}
\NormalTok{            relation, edus }\OperatorTok{=}\NormalTok{ token[}\StringTok{\textquotesingle{}misc\textquotesingle{}}\NormalTok{][}\StringTok{\textquotesingle{}Discourse\textquotesingle{}}\NormalTok{].split(}\StringTok{\textquotesingle{}:\textquotesingle{}}\NormalTok{)}
            
            \CommentTok{\# Tenta dividir a anotação da relação em duas partes}
            \ControlFlowTok{try}\NormalTok{:}
                
                \CommentTok{\# Divide na string \textquotesingle{}{-}\textgreater{}\textquotesingle{} e atribui a \textquotesingle{}source\textquotesingle{}}
                \CommentTok{\# e \textquotesingle{}target\textquotesingle{}, respectivamente.}
\NormalTok{                source, target }\OperatorTok{=}\NormalTok{ edus.split(}\StringTok{\textquotesingle{}{-}\textgreater{}\textquotesingle{}}\NormalTok{)}
                
                \CommentTok{\# Subtrai 1 de \textquotesingle{}source\textquotesingle{} e de \textquotesingle{}target\textquotesingle{}, porque os identificadores}
                \CommentTok{\# usados no corpus GUM não começam em zero, ao contrário dos}
                \CommentTok{\# nossos Spans do spaCy que correspondem às unidades elementares}
                \CommentTok{\# de discurso. Também converte os números em inteiros}
                \CommentTok{\# (originalmente eles são strings!).}
\NormalTok{                source, target }\OperatorTok{=} \BuiltInTok{int}\NormalTok{(source) }\OperatorTok{{-}} \DecValTok{1}\NormalTok{, }\BuiltInTok{int}\NormalTok{(target) }\OperatorTok{{-}} \DecValTok{1}
            
            \CommentTok{\# O nó raiz da árvore RST não terá alvo, o que levanta um}
            \CommentTok{\# ValueError, já que só existe um item.}
            \ControlFlowTok{except} \PreprocessorTok{ValueError}\NormalTok{:}
                
                \CommentTok{\# Atribui o primeiro item de \textquotesingle{}edus\textquotesingle{} a \textquotesingle{}source\textquotesingle{} e define}
                \CommentTok{\# \textquotesingle{}target\textquotesingle{} como None.}
\NormalTok{                source, target }\OperatorTok{=}\NormalTok{ edus[}\DecValTok{0}\NormalTok{], }\VariableTok{None}
                
                \CommentTok{\# Subtrai 1 de \textquotesingle{}source\textquotesingle{}, conforme explicado acima.}
\NormalTok{                source }\OperatorTok{=} \BuiltInTok{int}\NormalTok{(source) }\OperatorTok{{-}} \DecValTok{1} 
                
            \CommentTok{\# Compila a definição da relação em uma tripla e a acrescenta}
            \CommentTok{\# à lista \textquotesingle{}relations\textquotesingle{}.}
\NormalTok{            relations.append((relation, source, target))}
            
        \CommentTok{\# Verifica se o Token atual é seguido por um espaço. Quando não é o}
        \CommentTok{\# caso — por exemplo, no Token ao final de um TokenList —, essa}
        \CommentTok{\# informação está disponível na chave \textquotesingle{}misc\textquotesingle{}.}
        \ControlFlowTok{if}\NormalTok{ token[}\StringTok{\textquotesingle{}misc\textquotesingle{}}\NormalTok{] }\KeywordTok{is} \KeywordTok{not} \VariableTok{None} \KeywordTok{and} \StringTok{\textquotesingle{}SpaceAfter\textquotesingle{}} \KeywordTok{in}\NormalTok{ token[}\StringTok{\textquotesingle{}misc\textquotesingle{}}\NormalTok{]:}
            
            \CommentTok{\# Se a chave \textquotesingle{}misc\textquotesingle{} guarda um dicionário com a chave \textquotesingle{}SpaceAfter\textquotesingle{}}
            \CommentTok{\# com o valor \textquotesingle{}No\textquotesingle{}, executa o bloco abaixo}
            \ControlFlowTok{if}\NormalTok{ token[}\StringTok{\textquotesingle{}misc\textquotesingle{}}\NormalTok{][}\StringTok{\textquotesingle{}SpaceAfter\textquotesingle{}}\NormalTok{] }\OperatorTok{==} \StringTok{\textquotesingle{}No\textquotesingle{}}\NormalTok{:}
                
                \CommentTok{\# Acrescenta o valor booleano False à lista \textquotesingle{}spaces\textquotesingle{}.}
\NormalTok{                spaces.append(}\VariableTok{False}\NormalTok{)}
            
        \CommentTok{\# Se a chave \textquotesingle{}SpaceAfter\textquotesingle{} não for encontrada em \textquotesingle{}misc\textquotesingle{}, o token é}
        \CommentTok{\# seguido por um espaço.}
        \ControlFlowTok{else}\NormalTok{:}

            \CommentTok{\# Acrescenta True à lista de espaços}
\NormalTok{            spaces.append(}\VariableTok{True}\NormalTok{)}
        
        \CommentTok{\# Atualiza o contador ao terminar de processar um objeto Token,}
        \CommentTok{\# somando 1 ao seu valor.}
\NormalTok{        counter }\OperatorTok{+=} \DecValTok{1}
\end{Highlighting}
\end{Shaded}

O código acima é uma rotina clássica de ``garimpo''. Varrer um
dicionário estruturado na unha para separar todos os grãos essenciais
que, mais na frente, ajudarão a moldar um objeto \emph{Doc} limpo,
mantendo viva a bagagem de informações aprofundadas do documento
original.

Como você notou, processar arquivos baseados num esquema rico (como o
CoNLL-U) não admite aventureiros. Entender exatamente as engrenagens de
cada modelo (aqui, a forma como o corpus GUM acondiciona regras
linguísticas dentro da coluna livre \texttt{misc}) exige leitura do
manual, para garantir que as fatias não derretem no final do trajeto.

Na vida normal de programadores de PLN, injetar texto simples no funil
do \texttt{Language} object para que o spaCy entregue, do outro lado, um
arquivo rico em morfologia e marcações estruturais é rotineiro, algo que
batemos o martelo na Parte II.

Entretanto, esse atalho \textbf{não serve para a missão atual}. E o
porquê disso é elementar: caso deixássemos a inteligência do spaCy
fragmentar a nosso texto de forma nativa e automática (ou seja,
tokenizando do seu próprio jeito), correríamos o perigo real do spaCy
tomar liberdades divergentes da estrutura em ``pedras fundamentais''
exigidas pelo CoNLL-U e nosso meticuloso rastreio cair por terra.
Perderíamos a sincronia geométrica das anotações discursivas!

A via correta passa pela ``fabricação artesanal'' do objeto \emph{Doc}.
Carregaremos os construtores na unha a partir da módulo matricial
\texttt{tokens}, somando, é claro, a bagagem intelectual (o idioma de
vocabulário) presente num clássico modelo inglês simplificado (o
modesto, mas competente, \texttt{en\_core\_web\_sm}).

\begin{Shaded}
\begin{Highlighting}[]
\CommentTok{\# Importa a classe Doc e a biblioteca spaCy}
\ImportTok{from}\NormalTok{ spacy.tokens }\ImportTok{import}\NormalTok{ Doc}
\ImportTok{import}\NormalTok{ spacy}

\CommentTok{\# Carrega um modelo de linguagem pequeno para o inglês; armazena em \textquotesingle{}nlp\textquotesingle{}}
\NormalTok{nlp }\OperatorTok{=}\NormalTok{ spacy.load(}\StringTok{\textquotesingle{}en\_core\_web\_sm\textquotesingle{}}\NormalTok{)}
\end{Highlighting}
\end{Shaded}

Com o arcabouço pronto, injetaremos no objeto as matrizes que
laboriosamente resgatamos antes: a listagem das palavras
(\texttt{words}), seus indícios delimitadores interespaciais
(\texttt{spaces}) e os batizadores de largada frasal
(\texttt{sent\_starts}).

Para ilustrar o cenário atual, as impressões a seguir trazem as 15
``primeiras pedras'' agrupadas na montagem paralela das referidas
matrizes.

\begin{Shaded}
\begin{Highlighting}[]
\CommentTok{\# Percorre os 15 primeiros itens das listas \textquotesingle{}words\textquotesingle{}, \textquotesingle{}spaces\textquotesingle{} e \textquotesingle{}sent\_starts\textquotesingle{}.}
\CommentTok{\# Usa a função zip() para buscar os itens das listas simultaneamente.}
\ControlFlowTok{for}\NormalTok{ word, space, sent\_start }\KeywordTok{in} \BuiltInTok{zip}\NormalTok{(words[:}\DecValTok{15}\NormalTok{], spaces[:}\DecValTok{15}\NormalTok{], sent\_starts[:}\DecValTok{15}\NormalTok{]):}
    
    \CommentTok{\# Imprime o item atual de cada lista}
    \BuiltInTok{print}\NormalTok{(word, space, sent\_start)}
\end{Highlighting}
\end{Shaded}

\begin{verbatim}
How True True
to True False
Cope True False
With True False
a True False
Double True False
Parachute True False
Failure True False
While True True
skydiving False False
, True False
it True False
is True False
possible True False
( False False
\end{verbatim}

Além dessas três colunas vitais, é nossa tarefa amarrar essas diretrizes
ao vocabulário inerente do inglês, repassando o coração do modelo
(\texttt{nlp.vocab}) por meio da propriedade \texttt{vocab} presente no
\texttt{Doc}.

\begin{Shaded}
\begin{Highlighting}[]
\CommentTok{\# Cria um objeto Doc do spaCy "manualmente"; atribui à variável \textquotesingle{}doc\textquotesingle{}}
\NormalTok{doc }\OperatorTok{=}\NormalTok{ Doc(vocab}\OperatorTok{=}\NormalTok{nlp.vocab, }
\NormalTok{          words}\OperatorTok{=}\NormalTok{words, }
\NormalTok{          spaces}\OperatorTok{=}\NormalTok{spaces,}
\NormalTok{          sent\_starts}\OperatorTok{=}\NormalTok{sent\_starts}
\NormalTok{          )}
\end{Highlighting}
\end{Shaded}

Missão cumprida! Obtemos com sucesso a gênese de um objeto \emph{Doc}
totalmente sob nosso comando métrico e delineado pelas fronteiras
autênticas do formato CoNLL-U.

\begin{Shaded}
\begin{Highlighting}[]
\CommentTok{\# Recupera os Tokens até o índice 15 do objeto Doc}
\NormalTok{doc[:}\DecValTok{15}\NormalTok{]}
\end{Highlighting}
\end{Shaded}

\begin{verbatim}
How to Cope With a Double Parachute Failure While skydiving, it is possible (
\end{verbatim}

Os 15 grãos textuais encadeados de forma primorosa evidenciam o encaixe,
ditado perfeitamente sob batuta da listagem \texttt{words} em compasso
fiel a intermitência ditada pelos sinalizadores em \texttt{spaces}.

Sem falar nos delimitadores frasais! As sentenças perfeitamente
modeladas pela marcação original brilham isoladas sem a interferência
natural sintática do spaCy (isso vive alojado nas dependências relativas
da gaveta \texttt{sents} do novo objeto \texttt{doc}).

\begin{Shaded}
\begin{Highlighting}[]
\CommentTok{\# Recupera as cinco primeiras sentenças do objeto Doc}
\BuiltInTok{list}\NormalTok{(doc.sents)[:}\DecValTok{5}\NormalTok{]}
\end{Highlighting}
\end{Shaded}

\begin{verbatim}
[How to Cope With a Double Parachute Failure,
 While skydiving, it is possible (yet extremely unlikely) that both your primary and reserve parachutes will malfunction, leaving you with no method of reducing your velocity.,
 In the vast majority of cases, this will not occur;,
 nevertheless, in this event these coping strategies may assist.,
 Steps]
\end{verbatim}

Entretanto, por focarmos nas limitações espaciais descartando a riqueza
sintática englobada na estrutura bruta CoNLL-U original (apenas
absorvemos o texto e discurso), os atributos internos desses
\emph{Tokens} (como tipo da palavra) encontram-se desamparados (em
estado oco).

Duvida? Veja o deserto retornado numa tentativa de pinçar as raízes
morfológicas aprofundadas no campo inerente \texttt{tag\_} pertencente
ao primeiro elo da frase.

\begin{Shaded}
\begin{Highlighting}[]
\CommentTok{\# Obtém a etiqueta refinada de classe gramatical do Token no índice 0}
\NormalTok{doc[}\DecValTok{0}\NormalTok{].tag\_}
\end{Highlighting}
\end{Shaded}

\begin{verbatim}
''
\end{verbatim}

Branco nítido, inexistência provada. O dilema atual é: o spaCy condena
objetos empacotados ``do lado de fora'' em sua linha normal de atuação
do modelo \texttt{nlp}. Precisamos nós mesmos encaminhar os papéis e
protocolar as operações injetando \texttt{doc} nos guichês
especializados isolados (\emph{tagger}, \emph{parser}, etc) disponíveis
pela estrada rotulada em \texttt{pipeline}! O mapa já havia sido
discutido de praxe em pormenor em trecho anterior (Parte II).

Iremos transitar com as premissas em \emph{loop} engarrafando e
enriquecendo nosso arcabouço através dos guichês encadeados:

\begin{Shaded}
\begin{Highlighting}[]
\CommentTok{\# Percorre os pares nome / componente disponíveis no atributo \textquotesingle{}pipeline\textquotesingle{}}
\CommentTok{\# do objeto Language \textquotesingle{}nlp\textquotesingle{}.}
\ControlFlowTok{for}\NormalTok{ name, component }\KeywordTok{in}\NormalTok{ nlp.pipeline:}
    
    \CommentTok{\# Usa uma string formatada para imprimir o \textquotesingle{}name\textquotesingle{} do componente}
    \BuiltInTok{print}\NormalTok{(}\SpecialStringTok{f"Now applying component }\SpecialCharTok{\{}\NormalTok{name}\SpecialCharTok{\}}\SpecialStringTok{ ..."}\NormalTok{)}
    
    \CommentTok{\# Fornece o objeto Doc existente ao componente e armazena as anotações}
    \CommentTok{\# atualizadas na variável de mesmo nome (\textquotesingle{}doc\textquotesingle{}).}
\NormalTok{    doc }\OperatorTok{=}\NormalTok{ component(doc)}
\end{Highlighting}
\end{Shaded}

\begin{verbatim}
Now applying component tok2vec ...
Now applying component tagger ...
Now applying component parser ...
Now applying component attribute_ruler ...
Now applying component lemmatizer ...
Now applying component ner ...
\end{verbatim}

Terminado o expediente operário o objeto exibe suas aquisições ao
retornar o valor preenchido para a busca em \texttt{tag\_}:

\begin{Shaded}
\begin{Highlighting}[]
\CommentTok{\# Obtém a etiqueta refinada de classe gramatical do Token no índice 0}
\NormalTok{doc[}\DecValTok{0}\NormalTok{].tag\_}
\end{Highlighting}
\end{Shaded}

\begin{verbatim}
'WRB'
\end{verbatim}

Melhor ainda! Temos à disposição o aparato de análises finas exclusivas
derivadas, evidenciadas de antemão através dos construtos sintáticos
agregados em sintagmas sob o título de \texttt{noun\_chunks}.

\begin{Shaded}
\begin{Highlighting}[]
\CommentTok{\# Obtém os cinco primeiros sintagmas nominais do objeto Doc}
\BuiltInTok{list}\NormalTok{(doc.noun\_chunks)[:}\DecValTok{5}\NormalTok{]}
\end{Highlighting}
\end{Shaded}

\begin{verbatim}
[a Double Parachute Failure,
 it,
 both your primary and reserve parachutes,
 you,
 no method]
\end{verbatim}

O detalhe fundamental: os limites frasais pautados ``na unha'' (sob
demanda \texttt{sent\_starts}) foram preservados pelo sistema que
historicamente tem a tara em atropelar e ditar suas próprias regras
baseadas na árvore de sintaxe das palavras em vez da pontuação pura!

\begin{Shaded}
\begin{Highlighting}[]
\CommentTok{\# Obtém as cinco primeiras sentenças do objeto Doc}
\BuiltInTok{list}\NormalTok{(doc.sents)[:}\DecValTok{5}\NormalTok{]}
\end{Highlighting}
\end{Shaded}

\begin{verbatim}
[How to Cope With a Double Parachute Failure,
 While skydiving, it is possible (yet extremely unlikely) that both your primary and reserve parachutes will malfunction, leaving you with no method of reducing your velocity.,
 In the vast majority of cases, this will not occur;,
 nevertheless, in this event these coping strategies may assist.,
 Steps]
\end{verbatim}

Cenário pacificado! De posse integral sob manto protetivo da anotação,
caminhamos aos derradeiros ritos em infundir na raiz o viés da dimensão
analítica calcada e mapeada em nível discursivo oriunda dos garimpos
estipulados via \texttt{discourse\_units}, atados nas nuances em
\texttt{sent\_types} e conexões retóricas nas amarras
\texttt{relations}.

\hypertarget{amarrando-o-modo-da-orauxe7uxe3o-uxe0s-entranhas-do-spacy}{%
\subsubsection{Amarrando o modo da oração às entranhas do
spaCy}\label{amarrando-o-modo-da-orauxe7uxe3o-uxe0s-entranhas-do-spacy}}

Começaremos carimbando o estado de humor (o tom gramatical, em suma, o
modo) atado nas bases do conjunto armazenado \texttt{sent\_types}.

Categorias típicas (como traçamos anteriormente) ditam a tônica em
pautas relativas ao uso do tipo ``infinitivo'', vertentes calcadas para
``imperativo'' e balizas de cunho ``declarativo''.

Um mergulho rápido na listagem comprova o resgate das etiquetas.

\begin{Shaded}
\begin{Highlighting}[]
\CommentTok{\# Imprime os tipos de sentença}
\NormalTok{sent\_types[:}\DecValTok{10}\NormalTok{]}
\end{Highlighting}
\end{Shaded}

\begin{verbatim}
['inf', 'decl', 'decl', 'decl', 'frag', 'imp', 'decl', 'imp', 'imp', 'imp']
\end{verbatim}

Nossa missão dita agora ``grudar'' as correspondentes chancelas na
traseira das sentenças que moldam nosso modelo \emph{Doc}. Para nos
certificarmos perante risco de incompatibilidade, atestaremos em conduta
primária o nivelamento quantitativo entre o catálogo listado e os lotes
dispostos encarnando o grupo frasal da arquitetura atual.

\begin{Shaded}
\begin{Highlighting}[]
\CommentTok{\# Verifica o tamanho da informação de tipo de sentença e}
\CommentTok{\# o número de sentenças do objeto Doc.}
\BuiltInTok{len}\NormalTok{(sent\_types) }\OperatorTok{==} \BuiltInTok{len}\NormalTok{(}\BuiltInTok{list}\NormalTok{(doc.sents))}
\end{Highlighting}
\end{Shaded}

\begin{verbatim}
True
\end{verbatim}

Para injetar as peculiaridades balizadas nas regras semânticas na
infraestrutura robusta do spaCy a conduta ortodoxa determina alocar
``tags customizadas'' (os ``custom attributes'') no bojo interno
associado com as particularidades dos conjuntos de texto rotulados pelas
diretrizes da propriedade \emph{Span}, que congrega, agrupa e encarna as
delimitações englobando agrupamentos formados por parcelas unidas no
manto de matriz \texttt{Token} na raiz interna atada e fixada de
preceito pela arquitetura global \texttt{Doc}.

Ora, se pararmos e inspecionarmos pontualmente, um lote sob regência
vinculada na pasta interna a nível \texttt{sents} compõe intrinsecamente
um viés formatado de acordo com uma vestimenta autêntica sob carimbo
\texttt{Span}.

\begin{Shaded}
\begin{Highlighting}[]
\CommentTok{\# Imprime a primeira sentença do Doc e seu type()}
\BuiltInTok{list}\NormalTok{(doc.sents)[}\DecValTok{0}\NormalTok{], }\BuiltInTok{type}\NormalTok{(}\BuiltInTok{list}\NormalTok{(doc.sents)[}\DecValTok{0}\NormalTok{])}
\end{Highlighting}
\end{Shaded}

\begin{verbatim}
(How to Cope With a Double Parachute Failure, spacy.tokens.span.Span)
\end{verbatim}

Através da constatação e premissa referendada, a investida que estampa
nossa meta em injetar a taxonomia vinculativa aos atributos providos nas
sentenças requereria e apela pela convocação pontual (uso do artifício
``import'') associando as bibliotecas balizadoras para lidar com
\emph{SpanGroup} e o genérico formatado em \emph{Span}.

\begin{Shaded}
\begin{Highlighting}[]
\CommentTok{\# Importa as classes SpanGroup e Span do spaCy}
\ImportTok{from}\NormalTok{ spacy.tokens.span\_group }\ImportTok{import}\NormalTok{ SpanGroup}
\ImportTok{from}\NormalTok{ spacy.tokens }\ImportTok{import}\NormalTok{ Span}
\end{Highlighting}
\end{Shaded}

O pilar de arrimo na infraestrutura para expandir a fronteira
padronizada da engrenagem passa por encampar de antemão perante a via
oficial (a rotina central descrita extensamente por aqui em meados da
Parte II) um novo formato peculiar de cadastro atrelado na raiz do
artefato central para \texttt{Span}. A premissa no caso se resume em
carregar na mochila dessa entidade a chancela extra referenciando o tom
``humorístico/mood''.

Já na esfera englobada operante na denominação particular em torno do
utilitário abrigando viés de cunho \texttt{SpanGroup}, ele serve ao
singelo fim propiciando agrupamentos customizados voltados a gerenciar
caixas fechadas acopladas abrigando trechos de \texttt{Spans},
viabilizando ancorá-los na esteira do atributo pai de nome estipulado
por \texttt{spans} nas entranhas da estrutura principal em \texttt{Doc}.

Logo, fabricaremos nosso conjunto em \texttt{SpanGroup}, englobando os
segmentos providos originariamente a nível \emph{sents} alocados no
âmago interno.

A montagem base exigirá perante uso imperativo a determinação do modelo
base a ser operado (em \texttt{doc}), estipulando chancelas em favor do
selo norteador pelo quesito em \texttt{name}, por fim ancorando na lista
em \texttt{spans} a leva e amostragem coligida do ninho mãe que acolheu
as referidas sentenças!

\begin{Shaded}
\begin{Highlighting}[]
\CommentTok{\# Cria um SpanGroup a partir dos Spans contidos no objeto Doc \textquotesingle{}doc\textquotesingle{}, que}
\CommentTok{\# foi criado a partir das anotações CoNLL{-}U. Esses Spans correspondem às}
\CommentTok{\# sentenças, cujos limites definimos manualmente. Atribui à variável}
\CommentTok{\# \textquotesingle{}sent\_group\textquotesingle{}.}
\NormalTok{sent\_group }\OperatorTok{=}\NormalTok{ SpanGroup(doc}\OperatorTok{=}\NormalTok{doc, name}\OperatorTok{=}\StringTok{"sentences"}\NormalTok{, spans}\OperatorTok{=}\BuiltInTok{list}\NormalTok{(doc.sents))}
\end{Highlighting}
\end{Shaded}

Pronto. Deixamos as instâncias consolidadas a fim da inserção pontual,
alocando a rubrica em favor de fixar raízes perante os arranjos no
diretório em prol dos referidos sintagmas abrigados no domínio basilar
no diretório central:

\begin{Shaded}
\begin{Highlighting}[]
\CommentTok{\# Atribui o SpanGroup ao atributo \textquotesingle{}spans\textquotesingle{} do objeto Doc, sob a chave}
\CommentTok{\# \textquotesingle{}sentences\textquotesingle{}.}
\NormalTok{doc.spans[}\StringTok{\textquotesingle{}sentences\textquotesingle{}}\NormalTok{] }\OperatorTok{=}\NormalTok{ sent\_group}
\end{Highlighting}
\end{Shaded}

Com o reduto operante garantido, fixamos o registro estendendo a alçada
do novo selo ``mood'' (como propriedade de viés particular a ser fixado)
sob as regras do arcabouço \emph{Span}, inicializado de forma segura
assumindo um valor cego padronizado de segurança estatuído nulo (None).

\begin{Shaded}
\begin{Highlighting}[]
\CommentTok{\# Registra o atributo customizado \textquotesingle{}mood\textquotesingle{} na classe Span}
\NormalTok{Span.set\_extension(}\StringTok{\textquotesingle{}mood\textquotesingle{}}\NormalTok{, default}\OperatorTok{=}\VariableTok{None}\NormalTok{)}
\end{Highlighting}
\end{Shaded}

Nesse exato patamar é possível que pasmemos e apliquemos o fluxo
rotineiro iterativo encarnando e transpondo engate simultâneo pautado em
chaves pares casando ``cada tom gramatical com sua oração nativa''. O
Python cuida brilhantemente bem dessa proeza invocando sua magia
incorporada de modo inerente em pauta através do operador unificador com
viés emparelhador em ``zip()''.

Lembrando, não é mágica --- simplesmente vamos ditar a sintaxe base
associando os componentes alinhados atrelados à alocada e nomeada
\texttt{mood} amarrados por seu turno associando as frases batizadas
provisoriamente sob alcunha da vertente em \texttt{span}.

Fincamos chancelas englobando um detalhe sutil e clássico do spaCy na
transição de preenchimento balizando formatação na marca referenciada ao
componente via notação singular em sublinhado ``\texttt{\_}''. Trata-se
do reduto privativo reservado a inserções que partem pelo usuário
(detalhe já repassado anteriormente!).

\begin{Shaded}
\begin{Highlighting}[]
\CommentTok{\# Percorre os pares de itens da lista e do SpanGroup}
\ControlFlowTok{for}\NormalTok{ mood, span }\KeywordTok{in} \BuiltInTok{zip}\NormalTok{(sent\_types, doc.spans[}\StringTok{\textquotesingle{}sentences\textquotesingle{}}\NormalTok{]):}
    
    \CommentTok{\# Atribui o valor de \textquotesingle{}mood\textquotesingle{} ao atributo customizado de mesmo}
    \CommentTok{\# nome, pertencente ao objeto Span.}
\NormalTok{    span.\_.mood }\OperatorTok{=}\NormalTok{ mood}
\end{Highlighting}
\end{Shaded}

Agora nossos blocos textuais detêm uma carga informacional aprimorada!

Que tal pautar averiguação conferindo os arranjos nas chancelas
amparadas na sentença alojada no endereço do marcador de número 8?

\begin{Shaded}
\begin{Highlighting}[]
\CommentTok{\# Recupera a informação de modo verbal da sentença no índice 8}
\NormalTok{doc.spans[}\StringTok{\textquotesingle{}sentences\textquotesingle{}}\NormalTok{][}\DecValTok{8}\NormalTok{], doc.spans[}\StringTok{\textquotesingle{}sentences\textquotesingle{}}\NormalTok{][}\DecValTok{8}\NormalTok{].\_.mood}
\end{Highlighting}
\end{Shaded}

\begin{verbatim}
(If both your primary and reserve chutes have malfunctioned, signal immediately to a fellow jumper who has not yet deployed their chute, waving your arms and signalling that your own is not functional.,
 'imp')
\end{verbatim}

E aí está. O modelo acusa na exatidão, e o diagnóstico assinala que a
frase possui um modo ``imperativo'' enraizado em seu bojo gramatical.
Funcionou como um relógio!

\hypertarget{tecendo-o-discurso-com-as-anotauxe7uxf5es-das-relauxe7uxf5es-rst}{%
\subsubsection{Tecendo o discurso com as anotações das relações
RST}\label{tecendo-o-discurso-com-as-anotauxe7uxf5es-das-relauxe7uxf5es-rst}}

Chegou a hora e a vez de amarrarmos no pilar da ferramenta o estigma
mais profundo da análise. O limite exato estipulando chancelas perante
blocos de raciocínios independentes da teoria retórica se encontrava
enlatado desde outrora nos arranjos guardados sob posse no acervo com
estância referenciada nominalmente pela bagagem
\texttt{discourse\_units}.

\begin{Shaded}
\begin{Highlighting}[]
\CommentTok{\# Verifica os 10 primeiros itens da lista \textquotesingle{}discourse\_units\textquotesingle{}}
\NormalTok{discourse\_units[:}\DecValTok{10}\NormalTok{]}
\end{Highlighting}
\end{Shaded}

\begin{verbatim}
[0, 8, 11, 14, 19, 29, 34, 39, 51, 62]
\end{verbatim}

Os indícios aqui apontam friamente onde uma ``unidade retórica''
brotava! E como transpor a estatística num formato orgânico? O viés
lógico demanda pinçarmos recortes vivos no coração da estrutura original
em \emph{Doc} operando a faca exatamente por onde pautam as cercas
balizadoras indicadas a nível de índices (ex: corta de zero a oito, de
oito a onze, de onze a catorze, e sucessivamente até esvair o conteúdo)!

Esse desmembramento é facilmente vencido em pauta com o trato do
encadeamento oriundo por via sequencial de loops operados em fôlego em
\texttt{range()} fatiando pontualmente de praxe:

\begin{Shaded}
\begin{Highlighting}[]
\CommentTok{\# Cria uma lista vazia para guardar as fatias do objeto Doc que correspondem}
\CommentTok{\# às unidades de discurso.}
\NormalTok{edu\_spans }\OperatorTok{=}\NormalTok{ []}

\CommentTok{\# Percorre os limites das unidades de discurso usando a função range() do Python.}
\CommentTok{\# Isso nos dá números, que usamos para indexar a lista \textquotesingle{}discourse\_units\textquotesingle{}, a qual}
\CommentTok{\# contém os índices que marcam o início de cada unidade de discurso.}
\ControlFlowTok{for}\NormalTok{ i }\KeywordTok{in} \BuiltInTok{range}\NormalTok{(}\BuiltInTok{len}\NormalTok{(discourse\_units)):}
    
    \CommentTok{\# Tenta executar o bloco de código a seguir}
    \ControlFlowTok{try}\NormalTok{:}
        
        \CommentTok{\# Obtém o item atual da lista \textquotesingle{}discourse\_units\textquotesingle{} e o item seguinte; atribui}
        \CommentTok{\# às variáveis \textquotesingle{}start\textquotesingle{} e \textquotesingle{}end\textquotesingle{}.}
\NormalTok{        start, end }\OperatorTok{=}\NormalTok{ discourse\_units[i], discourse\_units[i }\OperatorTok{+} \DecValTok{1}\NormalTok{]}
    
    \CommentTok{\# Se o item seguinte não existir, porque chegamos ao último item da lista,}
    \CommentTok{\# um IndexError será levantado, e nós o capturamos aqui.}
    \ControlFlowTok{except} \PreprocessorTok{IndexError}\NormalTok{:}
        
        \CommentTok{\# Atribui o início da unidade de discurso normalmente e usa o tamanho do}
        \CommentTok{\# objeto Doc como valor de \textquotesingle{}end\textquotesingle{}, marcando o fim da unidade de discurso.}
\NormalTok{        start, end }\OperatorTok{=}\NormalTok{ discourse\_units[i], }\BuiltInTok{len}\NormalTok{(doc)}

    \CommentTok{\# Usa as variáveis \textquotesingle{}start\textquotesingle{} e \textquotesingle{}end\textquotesingle{} para fatiar o objeto Doc; acrescenta o}
    \CommentTok{\# objeto Span resultante à lista \textquotesingle{}edu\_spans\textquotesingle{}.}
\NormalTok{    edu\_spans.append(doc[start:end])}
\end{Highlighting}
\end{Shaded}

Com o esquadrinhamento selado na gaveta de repouso por trás das
fronteiras contíguas moldadas na matriz listada \texttt{edu\_spans},
temos a massa fundamental fragmentada.

\begin{Shaded}
\begin{Highlighting}[]
\CommentTok{\# Obtém os sete primeiros Spans da lista \textquotesingle{}edu\_spans\textquotesingle{}}
\NormalTok{edu\_spans[:}\DecValTok{7}\NormalTok{]}
\end{Highlighting}
\end{Shaded}

\begin{verbatim}
[How to Cope With a Double Parachute Failure,
 While skydiving,,
 it is possible,
 (yet extremely unlikely),
 that both your primary and reserve parachutes will malfunction,,
 leaving you with no method,
 of reducing your velocity.]
\end{verbatim}

É latente a total indiferença entre o que se estipula chancelas num
pilar da gramática ordinária para fixar sentenças longas em relação às
métricas para moldar as chancelas da unidade discursiva, descompasso
total!

Mas há peças cruciais pendentes focando os atributos necessários
amarrando a premissa de relação e subordinação de um componente sob a
mira do vizinho discursivo que moldará seu papel prático.

Logo criaremos mais três atributos sob estigma ``feito em casa'' no
esqueleto em \emph{Span}!

\begin{Shaded}
\begin{Highlighting}[]
\CommentTok{\# Registra três atributos customizados para objetos Span, correspondentes ao}
\CommentTok{\# id da unidade elementar de discurso, ao id do elemento que atua como alvo}
\CommentTok{\# e ao nome da relação.}
\NormalTok{Span.set\_extension(}\StringTok{\textquotesingle{}edu\_id\textquotesingle{}}\NormalTok{, default}\OperatorTok{=}\VariableTok{None}\NormalTok{)}
\NormalTok{Span.set\_extension(}\StringTok{\textquotesingle{}target\_id\textquotesingle{}}\NormalTok{, default}\OperatorTok{=}\VariableTok{None}\NormalTok{)}
\NormalTok{Span.set\_extension(}\StringTok{\textquotesingle{}relation\textquotesingle{}}\NormalTok{, default}\OperatorTok{=}\VariableTok{None}\NormalTok{)}
\end{Highlighting}
\end{Shaded}

Nesta engenharia: - \texttt{edu\_id}: Grava o selo de identificação
único para a porção elementar retórica. - \texttt{target\_id}: Guarda
quem é a outra porção conectada que sofre impacto prático. -
\texttt{relation}: Rotula a natureza que define a ponte psicológica a
entrelaçá-las.

E vamos novamente aninhar nossa ninhada fracionada de volta ao berço em
\emph{Doc} gerindo-as através da formatação de empacotamento com
maestria englobada da já conhecida classe \emph{SpanGroup}.

\begin{Shaded}
\begin{Highlighting}[]
\CommentTok{\# Cria um objeto SpanGroup a partir dos Spans da lista \textquotesingle{}edu\_spans\textquotesingle{}}
\NormalTok{edu\_group }\OperatorTok{=}\NormalTok{ SpanGroup(doc}\OperatorTok{=}\NormalTok{doc, name}\OperatorTok{=}\StringTok{"edus"}\NormalTok{, spans}\OperatorTok{=}\NormalTok{edu\_spans)}

\CommentTok{\# Atribui o SpanGroup sob a chave \textquotesingle{}edus\textquotesingle{}}
\NormalTok{doc.spans[}\StringTok{\textquotesingle{}edus\textquotesingle{}}\NormalTok{] }\OperatorTok{=}\NormalTok{ edu\_group}
\end{Highlighting}
\end{Shaded}

A etapa do arremate pede apenas que a gente alimente nossos ``tubos de
ensaio'' vazios com o verdadeiro sangue englobado na matriz inicial (a
tripla de tuplas guardada carinhosamente no banco com título de
\texttt{relations}).

Ao puxarmos a lista originária veremos o amparo provido calcado em
formatação atrelada a ditar ``quem se comunica e como se comunica''.

\begin{Shaded}
\begin{Highlighting}[]
\CommentTok{\# Imprime as cinco primeiras definições de relação da lista \textquotesingle{}relations\textquotesingle{}}
\NormalTok{relations[:}\DecValTok{5}\NormalTok{]}
\end{Highlighting}
\end{Shaded}

\begin{verbatim}
[('preparation', 0, 10),
 ('circumstance', 1, 2),
 ('background', 2, 8),
 ('concession', 3, 2),
 ('same-unit', 4, 2)]
\end{verbatim}

Resta a burocracia de encaminhar a transferência distribuindo a papelada
para a respectiva subseção criada (\texttt{edus} no guarda-chuva de
\texttt{spans}).

\begin{Shaded}
\begin{Highlighting}[]
\CommentTok{\# Percorre cada tripla da lista \textquotesingle{}relations\textquotesingle{}}
\ControlFlowTok{for}\NormalTok{ relation }\KeywordTok{in}\NormalTok{ relations:}
    
    \CommentTok{\# Desempacota a tripla em três variáveis}
\NormalTok{    rel\_name, source, target }\OperatorTok{=}\NormalTok{ relation[}\DecValTok{0}\NormalTok{], relation[}\DecValTok{1}\NormalTok{], relation[}\DecValTok{2}\NormalTok{]}
    
    \CommentTok{\# Usa o identificador em \textquotesingle{}source\textquotesingle{} para indexar os objetos Span sob a}
    \CommentTok{\# chave \textquotesingle{}edus\textquotesingle{}. Em seguida, acessa os atributos customizados e define os}
    \CommentTok{\# valores de \textquotesingle{}edu\_id\textquotesingle{}, \textquotesingle{}target\_id\textquotesingle{} e \textquotesingle{}relation\textquotesingle{}.}
\NormalTok{    doc.spans[}\StringTok{\textquotesingle{}edus\textquotesingle{}}\NormalTok{][source].\_.edu\_id }\OperatorTok{=}\NormalTok{ source}
\NormalTok{    doc.spans[}\StringTok{\textquotesingle{}edus\textquotesingle{}}\NormalTok{][source].\_.target\_id }\OperatorTok{=}\NormalTok{ target}
\NormalTok{    doc.spans[}\StringTok{\textquotesingle{}edus\textquotesingle{}}\NormalTok{][source].\_.relation }\OperatorTok{=}\NormalTok{ rel\_name}
\end{Highlighting}
\end{Shaded}

Bingo! A trama retórica se encontra solidamente fixada ao osso da
engrenagem!

Vamos averigar o status retórico encravado no pedaço catalogado na
posição sob índice de marcação de número 1:

\begin{Shaded}
\begin{Highlighting}[]
\CommentTok{\# Recupera o atributo customizado \textquotesingle{}relation\textquotesingle{} do Span no índice 1}
\NormalTok{doc.spans[}\StringTok{\textquotesingle{}edus\textquotesingle{}}\NormalTok{][}\DecValTok{1}\NormalTok{].\_.relation}
\end{Highlighting}
\end{Shaded}

\begin{verbatim}
'circumstance'
\end{verbatim}

Ele acusa operar na pauta ``circunstancial'' ditando o viés dependendo
de amarras de praxe atreladas ao reduto apontando rumo na outra oração
da extremidade espelhada balizada ao registro com ``target\_id''. Mas
quem é o comparsa alvejado na mira do ``target\_id''?

\begin{Shaded}
\begin{Highlighting}[]
\CommentTok{\# Obtém o Span no índice 1 e o Span referenciado por seu atributo \textquotesingle{}target\_id\textquotesingle{}}
\NormalTok{doc.spans[}\StringTok{\textquotesingle{}edus\textquotesingle{}}\NormalTok{][}\DecValTok{1}\NormalTok{], doc.spans[}\StringTok{\textquotesingle{}edus\textquotesingle{}}\NormalTok{][doc.spans[}\StringTok{\textquotesingle{}edus\textquotesingle{}}\NormalTok{][}\DecValTok{1}\NormalTok{].\_.target\_id]}
\end{Highlighting}
\end{Shaded}

\begin{verbatim}
(While skydiving,, it is possible)
\end{verbatim}

Eis os gêmeos siameses desnudados: As porções encadeadas comprovando
empiricamente e validando as amarras estipuladas entre as unidades
elementares discursivas atreladas sob preceito referenciando chancelas
do tipo que dita uma circunstância.

\hypertarget{atalhos-para-importar-formatos-conll-u-para-dentro-do-spacy}{%
\subsection{Atalhos para importar formatos CoNLL-U para dentro do
spaCy}\label{atalhos-para-importar-formatos-conll-u-para-dentro-do-spacy}}

Caso não haja o apego emocional no ofício minucioso ao querer incorporar
customizações artesanais para inflar a estância do modelo atrelada ao
\texttt{Doc} e a intenção recaia friamente num interesse mecânico de
fazer um passe mágico com viés restrito na pauta focada convertendo
CoNLL-U num repositório puro formatado no padrão referenciando ao objeto
balizado em spaCy, uma luz brilhante te aguarda: o método atalho sob o
nome em utilitário de plantão chamado \texttt{conllu\_to\_docs()}.

Ele foi cunhado para viver no nicho englobando o amparo em setores base
de processamentos para \texttt{training}, que figura na alçada balizando
modelos alimentadores no treinamento focado sob as artérias de bases
encorpadas pautando aprendizado computacional.

\begin{Shaded}
\begin{Highlighting}[]
\CommentTok{\# Importa a função \textquotesingle{}conllu\_to\_docs\textquotesingle{} e a classe Doc}
\ImportTok{from}\NormalTok{ spacy.training.converters }\ImportTok{import}\NormalTok{ conllu\_to\_docs}
\ImportTok{from}\NormalTok{ spacy.tokens }\ImportTok{import}\NormalTok{ Doc}
\end{Highlighting}
\end{Shaded}

Este artefato ``mastiga'' as entradas fornecidas via formato basilar
operado na vertente \texttt{string}.

Empregamos no momento a massa bruta \texttt{annotations} fornecendo-a à
máquina. E para que ela opere ``quietinha'', escoramos o amparo
chancelando premissas operantes na vertente que manda ela trancar a voz
no quesito da opção atada e estatuída na marca \texttt{no\_print}
definida num ``ok'' positivo.

A rotina irá cuspir a produção (sob formato Pythonic em `generator'), de
modo a requerer do usuário providências para converter tudo à moda
clássica na moldura para base em formato ``lista'' (list).

\begin{Shaded}
\begin{Highlighting}[]
\CommentTok{\# Fornece a string em \textquotesingle{}annotations\textquotesingle{} à função \textquotesingle{}conllu\_to\_docs\textquotesingle{}.}
\CommentTok{\# Define \textquotesingle{}no\_print\textquotesingle{} como True e converte o resultado em lista; armazena em \textquotesingle{}docs\textquotesingle{}.}
\NormalTok{docs }\OperatorTok{=} \BuiltInTok{list}\NormalTok{(conllu\_to\_docs(annotations, no\_print}\OperatorTok{=}\VariableTok{True}\NormalTok{))}

\CommentTok{\# Obtém os dois primeiros itens da lista resultante}
\NormalTok{docs[:}\DecValTok{2}\NormalTok{]}
\end{Highlighting}
\end{Shaded}

\begin{verbatim}
[How to Cope With a Double Parachute Failure While skydiving, it is possible (yet extremely unlikely) that both your primary and reserve parachutes will malfunction, leaving you with no method of reducing your velocity. In the vast majority of cases, this will not occur; nevertheless, in this event these coping strategies may assist. Steps Remain calm. This may seem obvious, but deep, even breathing (despite the rushing slipstream) and controlling your heart-rate are essential to your continued survival. Entreat assistance. If both your primary and reserve chutes have malfunctioned, signal immediately to a fellow jumper who has not yet deployed their chute, waving your arms and signalling that your own is not functional. If this is a solo jump, then skip to step 6. ,
 Prepare for deployment. After linking arms with your fellow jumper, you will need to hook your arms through their chest strap, or through both sides of the front of their harness, as far as you can, then grab onto your own strap. Deploy. The shock of the chute deployment will be intense. The G-forces will multiply your body weight, making it impossible to hold on; this is why hooking your arms through the harness is essential. It is likely that the shock will dislocate or break both your arms; nevertheless this is a small price to pay for your life. Prepare for impact. If your companion's canopy has successfully opened, then both your chances of survival have dramatically increased. ]
\end{verbatim}

Pronto, criamos as fatias perfeitamente moldadas ao \texttt{Doc}. O
pequeno tropeço na engrenagem que assombra essa facilidade balizada em
padrão de fábrica: ela quebra arbitrariamente os documentos em maços
picotados agrupando as ocorrências pautando-se sob limites focados ao
amparo englobando 10 em 10 sentenças em caixinhas alocadas independentes
de \texttt{Doc}s individuais (um tanto desastroso).

Uma fusão coerente focada na reparação recai operando base numa vertente
ditada a promover aglutinações embasadas a recompor o texto global num
arcabouço atado na pauta unitária no viés em amarrações ancoradas nas
chamadas à subrotina originada do artifício balizador estatuído sob a
função \texttt{from\_docs()} provinda nativa na premissa com chancela
perante molde embasando no coração \emph{Doc}.

\begin{Shaded}
\begin{Highlighting}[]
\CommentTok{\# Combina os objetos Doc da lista \textquotesingle{}docs\textquotesingle{} em um único Doc; atribui a \textquotesingle{}doc\textquotesingle{}}
\NormalTok{doc }\OperatorTok{=}\NormalTok{ Doc.from\_docs(docs)}

\CommentTok{\# Verifica o tipo e o tamanho da variável}
\BuiltInTok{type}\NormalTok{(doc), }\BuiltInTok{len}\NormalTok{(doc)}
\end{Highlighting}
\end{Shaded}

\begin{verbatim}
(spacy.tokens.doc.Doc, 890)
\end{verbatim}

Milagre operado: Um belo e reluzente documento singular composto
majestosamente em 890 matrizes providas de \emph{Tokens} independentes!
E é claro, com todas as virtudes de informações linguísticas incrustadas
de nascença!

\begin{Shaded}
\begin{Highlighting}[]
\CommentTok{\# Percorre os 8 primeiros Tokens usando a função range()}
\ControlFlowTok{for}\NormalTok{ token\_ix }\KeywordTok{in} \BuiltInTok{range}\NormalTok{(}\DecValTok{0}\NormalTok{, }\DecValTok{8}\NormalTok{):}
    
    \CommentTok{\# Usa o número atual em \textquotesingle{}token\_ix\textquotesingle{} para buscar um Token no Doc.}
    \CommentTok{\# Atribui o objeto Token à variável \textquotesingle{}token\textquotesingle{}.}
\NormalTok{    token }\OperatorTok{=}\NormalTok{ doc[token\_ix]}
    
    \CommentTok{\# Imprime o Token e suas anotações linguísticas}
    \BuiltInTok{print}\NormalTok{(token, token.tag\_, token.pos\_, token.morph, token.dep\_, token.head)}
\end{Highlighting}
\end{Shaded}

\begin{verbatim}
How WRB SCONJ PronType=Int mark Cope
to TO PART  mark Cope
Cope VB VERB VerbForm=Inf ROOT Cope
With IN ADP  case Failure
a DT DET Definite=Ind|PronType=Art det Failure
Double JJ ADJ Degree=Pos amod Failure
Parachute NN NOUN Number=Sing compound Failure
Failure NN NOUN Number=Sing obl Cope
\end{verbatim}

Entretanto, não cante vitória de cara perante as virtudes. Ao esbarrar
na tentação em requerer e focar em propriedades estatuindo sintagmas
compostos enquadrando alocações na premissa originária de entidades em
teor \emph{noun\_chunks}, a interface irá corar de vergonha na sua tela
gritando e acusando por chancelas atadas em lacunas de dados que
causarão a fúria em erros (``NotImplementedError'').

\begin{Shaded}
\begin{Highlighting}[]
\CommentTok{\# Obtém os sintagmas nominais do Doc}
\BuiltInTok{list}\NormalTok{(doc.noun\_chunks)}
\end{Highlighting}
\end{Shaded}

\begin{verbatim}
---------------------------------------------------------------------------
NotImplementedError                       Traceback (most recent call last)
Cell In[47], line 2
      1 # Get the noun chunks in the Doc
----> 2 list(doc.noun_chunks)

File ~/nlp/lib/python3.9/site-packages/spacy/tokens/doc.pyx:853, in noun_chunks()

NotImplementedError: [E894] The 'noun_chunks' syntax iterator is not implemented for language ''.
\end{verbatim}

O revés ataca em formato oriundo puramente operando a nível estrito: o
\texttt{Doc} gestado não porta qualquer cordão umbilical ligando a
natureza a modelos estatuídos operantes do nível matricial de
\textbf{Language} (\emph{Vocabulary} e outros parâmetros que determinam
o que a máquina sabe do idioma). E os extratos ``noun chunks'' em
essência sugam de praxe, obrigatoriamente, embasamentos providos nestes
dicionários linguísticos acoplados a suas raízes que chancelam sintaxe
perante bases que o sistema ignora no instante atual.

Para pautar chancelas de praxe de cura e reabilitação não dispomos da
possibilidade provida via injeção operando rotinas à força de maneira
puramente e simplesmente ditar as cartilhas linguísticas à mão. O
procedimento operante foca em balizas de ``cirurgias'' no armazenamento
do formato ``disco'' acionando base da estância chamada de pilar
\emph{DocBin}, abordada em aprofundamento pretérito provido perante o
horizonte da Parte II.

O pacote englobando as amarrações delineando no esteio pautado na pauta
oriunda a estância abarcando propriedades base \texttt{DocBin} ostenta
as prerrogativas chancelando enquadramento estipulando chancelas
voltadas a transpor dados focados à exportação externa da estrutura
pronta.

\begin{Shaded}
\begin{Highlighting}[]
\CommentTok{\# Importa o objeto DocBin do submódulo \textquotesingle{}tokens\textquotesingle{}}
\ImportTok{from}\NormalTok{ spacy.tokens }\ImportTok{import}\NormalTok{ DocBin}

\CommentTok{\# Cria um objeto DocBin vazio}
\NormalTok{doc\_bin }\OperatorTok{=}\NormalTok{ DocBin()}

\CommentTok{\# Adiciona o Doc atual ao DocBin}
\NormalTok{doc\_bin.add(doc)}
\end{Highlighting}
\end{Shaded}

Nesta pauta o revés reside ao invés do despejo em HD de arquivo bruto
nós faremos o saque na fenda operando o recuo, onde engajaremos o saque
imediato de volta para a pauta, puxando à luz do dia os extratos usando
artifícios da técnica invocada em subrotina \texttt{get\_docs()}, só que
agora empurramos pela guelra da máquina o pacote que traz a carga focada
em estatuir a língua que baliza o sistema de \texttt{Vocabulary}. O
artefato que devolve vida inteligente perante amarrações amparadas e
englobando o idioma no spaCy.

O fruto dessa manobra devolve uma safra ``generator'' e, logo, voltamos
à lida transformando isso no velho enquadramento em lista
(\texttt{list}).

\begin{Shaded}
\begin{Highlighting}[]
\CommentTok{\# Usa o método \textquotesingle{}get\_docs\textquotesingle{} para recuperar os Docs do DocBin}
\NormalTok{docs }\OperatorTok{=} \BuiltInTok{list}\NormalTok{(doc\_bin.get\_docs(vocab}\OperatorTok{=}\NormalTok{nlp.vocab))}
\end{Highlighting}
\end{Shaded}

Problema sanado, o acesso volta à tona operando de maneira brilhante! A
malha atada em viés referencial de propriedades contidas e embutidas sob
chancela nos redutos amparando embates voltados ao uso da estrutura
originária ditada por \emph{noun chunks} se submete e jorra farta sem o
engasgo da interdição perante a falta do idioma.

\begin{Shaded}
\begin{Highlighting}[]
\BuiltInTok{list}\NormalTok{(docs[}\DecValTok{0}\NormalTok{].noun\_chunks)[:}\DecValTok{10}\NormalTok{]}
\end{Highlighting}
\end{Shaded}

\begin{verbatim}
[it,
 and reserve parachutes,
 this,
 these coping strategies,
 Steps,
 This,
 both your primary and reserve chutes,
 who,
 this,
 you]
\end{verbatim}

Concluímos assim o mapeamento englobando o panorama referencial
pautando-se sob chancelas a pavimentar vias da sua trilha em lidar e
aprimorar a exploração das profundezas estruturais ancoradas no trato
por anotações abrigadas perante o estigma ditado por regras estatuídas
através da maestria provida nativa no formato operante sob regência do
esquema abarcado por CoNLL-U operando a nível estrito focando suas
intersecções atadas nas malhas do ecossistema e viés originário em
componentes providos ao amparo de funcionalidades estipulando
enquadramento balizado na estância de spaCy.

\end{document}
